\section{Introduction}
\label{sec:intro}

\subsection{Two objects that turn out to be one}

Start with a picture of a brain, or at least a cartoon of one. A population of
neurons responds to the things you see. For any particular thing --- a face, a
cup, your grandmother --- some of those neurons fire and the rest stay quiet.
Write down a $1$ for each active neuron and a $0$ for each silent one, and the
thing you were looking at has become a string of bits. Do this for many things
and you have a table: rows are the things, columns are the neurons, and each
cell says whether that neuron fired for that thing. This table is called a
\emph{neural code}, and in the biologically realistic case it is
\emph{sparse}: only a small fraction of the neurons are active for any one
stimulus \citep{barlow1972,foldiak1990,olshausen1996}.

Now notice something about the table that has nothing to do with biology. Take
two things --- say, a spaniel and a poodle --- and look at the neurons that
fire for both. If a large set of neurons is shared, the two things have
something in common. If the set of neurons active for ``spaniel'' happens to
\emph{contain} every neuron active for ``dog in general'', that containment is
a statement about categories: a spaniel is a dog. No wire in the brain has to
represent the arrow from spaniel to dog. The arrow is already there, implicit
in which codewords overlap which.

That is the observation \odag{} was built on \citep{foldiak2003,ontodag}:
\emph{a knowledge graph can be latent in a code}. The overlaps between
codewords are a category structure, and if you can compute with the overlaps
you can compute with the categories.

The observation has a mathematical home. It is called \textbf{Formal Concept
Analysis} (\fca{}), it was founded by Rudolf Wille in 1982
\citep{wille1982,ganter1999}, and it is precisely the theory of turning a
binary table into a hierarchy of concepts --- and the hierarchy back into the
table. \fca{}'s name for the table is a \emph{formal context}; its name for
the hierarchy is a \emph{concept lattice}. What \fca{} contributes is not a
vague analogy but an exact, invertible correspondence with forty years of
theory attached: algorithms, complexity results, a theory of implications
between attributes, an interactive protocol for growing a knowledge base with
an expert, and a well-understood way of handling numbers and intervals.

\begin{keyidea}
A neural code is a binary matrix. A binary matrix is a formal context. A
formal context determines a lattice of concepts. So a code determines a
concept hierarchy --- a \emph{mind} object recoverable from a \emph{brain}
object, and vice versa. This is the bridge, and \fca{} is the mathematics of
crossing it \citep{endres2008}.
\end{keyidea}

\odag{} is a working system that lives near this bridge but does not stand on
it. It is a running piece of software --- a category manager with a
command-line tool, a web interface, a Python API, and a persistence layer that
publishes to a decentralised storage network --- and its design is shaped by
engineering constraints that pure \fca{} never had to face: it must accept
edits one at a time, it must merge the work of writers who never talk to each
other, and its stored form must hash to the same 32 bytes whenever two people
know the same thing. Those constraints pushed it away from \fca{} in specific,
principled ways. This tutorial is about the resulting family resemblance: what
is shared, what differs, what each side should learn from the other.

\subsection{What this document is for}

Four questions, in order.

\begin{enumerate}[label=\textbf{Q\arabic*.},leftmargin=2.6em]
\item \textbf{How exactly do \odag{} and \fca{} correspond?} Not
  ``both are about hierarchies'' --- an exact dictionary, with a theorem
  (Section~\ref{sec:dictionary}).
\item \textbf{What can \odag{} take from \fca{}'s theory?} \fca{} is rich and
  mature. \odag{} is small and young. There is a shopping list, and some of it
  is affordable today (Section~\ref{sec:learn}).
\item \textbf{What must \odag{} \emph{not} adopt?} Some of \fca{}'s central
  commitments would destroy the properties that make \odag{} useful. Knowing
  which, and why, is more valuable than the shopping list
  (Section~\ref{sec:keep}).
\item \textbf{Who benefits?} Two constituencies, in detail: AI agents
  (Section~\ref{sec:agents}) and cryptographic / decentralised systems
  (Section~\ref{sec:crypto}). The claim is that they need the \emph{same} four
  properties, and that those properties are exactly what \odag{}'s refusals
  buy.
\end{enumerate}

The document is written as a tutorial. It assumes no order theory, no lattice
theory, no \fca{}, and no cryptography beyond ``a hash function turns data
into a short fingerprint''. Everything else is built up from school
mathematics: sets, subsets, and arrows between things. Where a technical term
is unavoidable it is defined at first use, boxed, and repeated in the glossary
(Appendix~\ref{sec:glossary}).

\subsection{One example, used throughout}
\label{sec:running-example}

To keep the two formalisms comparable, one small example runs through the
whole paper. You are planning a trip to Japan and you have six documents.

\begin{center}
\small
\begin{tabular}{lp{9.4cm}}
\toprule
\textbf{Document} & \textbf{What it is}\\
\midrule
\texttt{outbound.pdf}  & the flight out; booked; refundable\\
\texttt{inbound.pdf}   & the flight back; booked; not refundable\\
\texttt{kyoto-inn.pdf} & a hotel in Kyoto; booked; refundable\\
\texttt{rail-pass.pdf} & a Japan Rail pass; booked; refundable\\
\texttt{paris-hop.pdf} & an unrelated flight to Paris; booked\\
\texttt{visa-note.md}  & notes on the visa; not a booking at all\\
\bottomrule
\end{tabular}
\end{center}

There are five properties in play: \textsf{Flight}, \textsf{Hotel},
\textsf{Japan}, \textsf{Booked}, \textsf{Refundable}. This is deliberately
tiny --- small enough that every computation in Parts~I and~II can be checked
by hand, and small enough that the concept lattice fits on a page. It is also
deliberately awkward in the right way: \texttt{outbound.pdf} is a flight
\emph{and} a Japan document \emph{and} refundable, which is exactly the
situation that a folder tree cannot represent and that motivates both
formalisms.

\subsection{The thread that connects the parts}

There is a single idea underneath all four questions, and it is worth stating
before the machinery arrives, because everything later is a variation on it.

Consider what it means for a system to \emph{know} that a spaniel is a dog.
There are two ways to hold that knowledge:

\begin{itemize}[leftmargin=1.4em]
\item \textbf{Say it.} Store an arrow from \textsf{spaniel} to \textsf{dog}.
  The knowledge is an \emph{assertion} --- somebody put it there, and it is
  there because they did.
\item \textbf{Derive it.} Store the facts about individual spaniels and
  individual dogs, and let the containment of one set of properties in another
  \emph{produce} the arrow. The knowledge is a \emph{consequence} --- nobody
  put it there; it follows.
\end{itemize}

\fca{} is the theory of the second option taken to its limit: given the data,
compute \emph{every} concept the data supports, and every relationship between
them. \odag{} is a system built firmly on the first: store what was asserted,
in the smallest form that implies the same things, and compute the rest at
query time. Nearly every difference in this paper --- in cost, in canonical
form, in mergeability, in whether concepts have stable names --- flows from
that one fork.

Neither choice is wrong. They answer different questions. \fca{} asks ``what
concepts does this data \emph{contain}?'' and gives a complete, closed,
finished answer about a fixed table. \odag{} asks ``what have people said, and
what follows from it right now?'' and must give an answer that survives being
edited, merged with a stranger's answer, and fingerprinted for a blockchain.

\begin{figure}[t]
\centering
\begin{tikzpicture}[
  font=\small,
  box/.style={draw, rounded corners=3pt, minimum height=1.05cm,
              text width=2.85cm, align=center, inner sep=4pt},
  lbl/.style={font=\scriptsize\itshape, text=odgrey, align=center},
  ar/.style={-{Stealth[length=2.4mm]}, thick}]

\node[box, fill=odsand] (code) at (0,0)
  {\textbf{A code}\\[1pt]\scriptsize binary matrix\\\scriptsize
   things $\times$ features};
\node[box, fill=white, draw=odblue] (ctx) at (5.0,0)
  {\textbf{Formal context}\\[1pt]\scriptsize $\ctx=(G,M,I)$\\\scriptsize
   \fca{}'s input};
\node[box, fill=white, draw=odblue] (lat) at (10.0,0)
  {\textbf{Concept lattice}\\[1pt]\scriptsize $\Bcl(\ctx)$\\\scriptsize
   every closed set};
\node[box, fill=white, draw=odteal] (dag) at (10.0,-2.6)
  {\textbf{\odag{}}\\[1pt]\scriptsize asserted edges only\\\scriptsize
   transitively reduced};
\node[box, fill=white, draw=odrust] (root) at (5.0,-2.6)
  {\textbf{A 32-byte root}\\[1pt]\scriptsize content address\\\scriptsize
   of the whole graph};
\node[box, fill=odsand] (uses) at (0,-2.6)
  {\textbf{Agents \& chains}\\[1pt]\scriptsize verify, cite,\\\scriptsize
   merge, settle};

\draw[ar] (code) -- node[lbl, above]{same\\object} (ctx);
\draw[ar] (ctx)  -- node[lbl, above]{close\\everything} (lat);
\draw[ar, dashed, odgrey] (lat) -- node[lbl, right, xshift=1mm]
  {keep only what\\was asserted\\(\S\ref{sec:dictionary})} (dag);
\draw[ar] (dag)  -- node[lbl, above]{canonical\\form + hash} (root);
\draw[ar] (root) -- node[lbl, above]{one address\\per meaning} (uses);
\end{tikzpicture}
\caption{The road this paper travels. A neural code and a formal context are
the same mathematical object seen by two communities. Closing the context
gives the concept lattice, the object \fca{} studies. \odag{} sits below it:
the asserted fragment, kept in a unique minimal form, which is what allows the
whole graph to have a single content address --- and that address is what AI
agents and cryptographic systems actually want.}
\label{fig:roadmap}
\end{figure}

Figure~\ref{fig:roadmap} shows the route. We will walk it from left to right.

\subsection{Reproducibility}

Every quantitative claim in this paper --- concept counts, lattice sizes,
canonical root hashes, the round-trip between codes and graphs --- is computed
by \file{paper/experiments/fca_experiments.py}, which uses only the Python
standard library plus \odag{} itself. Its full output is reproduced in
Appendix~\ref{sec:experiments}, and the figures drawn from data are generated
from the same run. Where a number appears in the text it is marked with a
box like this one:

\begin{measured}
Running example: 6 documents $\times$ 5 properties, $57\%$ of cells filled.
\fca{} finds \textbf{10} formal concepts (out of $2^5=32$ conceivable
property sets). The same facts as an \odag{}: \textbf{11} nodes and
\textbf{17} asserted edges.
\end{measured}
